\majorbreak
\section{Reduce generated proof repetition without changing the circuits}
\label{sec:circuits}

The largest source concentration is the pair \code{SplitCircuitSteps59}
and \code{GapCircuitSteps66}: 6,567 physical lines, comprising 999 gate
wrappers and 193 power lemmas. This is a strong source-maintenance target,
but it does not mean 1,192 independent mathematical ideas need to be replaced.
The gate family already consists of instances of one generic theorem.

\subsection{Use the semantic infrastructure already present}

The project has the key reusable results now. A proposal to introduce its
first generic DAG soundness or degree checker would miss these interfaces:

\begin{description}
\item[Gate and register equations]
\code{run_node_equation}, followed by \code{run_eq_of_node_equations}
and \code{value_eq_of_node_equations}, at
\src{ArithmeticDAGSemantics.lean}{58}{116}.
\item[Power labels and circuit rewrites]
\code{run_power} at \src{ArithmeticDAGPowers.lean}{24}{42}, and
\code{value_eq_of_powerRewrite} at
\src{ArithmeticDAGPowers.lean}{86}{131}.
\item[Degree certificates]
\code{DegreeCertificate} at \src{ArithmeticDAGDegree.lean}{21}{23},
with \code{value_degree_le} and \code{polynomial_degree_le} at
\src{ArithmeticDAGDegree.lean}{65}{90}.
\end{description}

\begin{table}[h]
\small
\begin{tabularx}{\linewidth}{@{}Yrrr@{}}
\toprule
Existing generated module & Total lines & Gate wrappers & Power lemmas\\
\midrule
\code{SplitCircuitSteps59} & 3,445 & 515 & 107\\
\code{GapCircuitSteps66} & 3,122 & 484 & 86\\
\bottomrule
\end{tabularx}
\caption{Gate spans are Split lines 16--2590 and Gap lines 16--2435.
Power spans are Split lines 2591--3444 and Gap lines 2436--3121.}
\end{table}

The measured gate spans occupy 4,995 lines and the power spans 1,540,
including blank lines. Together they cover 6,535 of the 6,567 lines.
These figures locate the repetition; they are not a deletion budget, because
clients currently use the generated names.

\subsection{Consolidate power proofs while retaining every public theorem}
\label{sec:powercert}

\textbf{Priority 2 and the safest circuit experiment.} Preserve the gate
wrappers, every public theorem header, and the literal DAG. Replace the
individual pure-power \code{ring} proofs with a finite exponent-label
certificate and the existing \code{run_power} theorem.

The current \code{GapCircuitSteps66.natural_power_25} at
\src{GapCircuitSteps66.lean}{2442}{2448} is:

\begin{Lean}
theorem natural_power_25 (p : Fin 11 → R) (x : Fin 9 → R) :
    naturalRegister p x 25 = naturalRegister p x 20 ^ 4 := by
  rw [natural_step_25]
  simp only [natural_power_24]
  ring
\end{Lean}

After adding the shared private certificate and generic specialization, the
candidate keeps that public proposition and changes only its proof:

\begin{Lean}
theorem natural_power_25 (p : Fin 11 → R) (x : Fin 9 → R) :
    naturalRegister p x 25 = naturalRegister p x 20 ^ 4 :=
  natural_power_all p x ⟨25, by decide⟩ (by decide +kernel)
\end{Lean}

The label table uses zero to mean \emph{not certified as a power}, not
exponent zero. The base has label one: register 36 in Split, register 20 in
Gap. Every positive non-base label belongs to a multiplication node whose
positive operand labels add to the result label. The existing topology
argument supports induction over register number, uniformly for every
\code{CommRing}.

For example, the proposed closed Gap obligation is:

\begin{Lean}
private theorem naturalPowerSteps : ∀ i : Fin 484,
    i.val ≠ 20 → naturalPowerLabel i.val ≠ 0 →
      PowerStep GapCircuit66.naturalNodes naturalPowerLabel i.val := by
  decide +kernel
\end{Lean}

The full \code{natural_power_all} proof invokes \code{run_power} and this
certificate. The archive contains the full label tables and complete
API-preserving candidate modules, not just the abbreviated listing here.

\begin{table}[h]
\small
\begin{tabularx}{\linewidth}{@{}Yrrr@{}}
\toprule
Final candidate & Physical lines & UTF-8 bytes & Per-power \code{ring} calls\\
\midrule
Split before & 3,445 & 169,808 & 107\\
Split candidate & 3,074 & 171,042 & 0\\
Gap before & 3,122 & 153,802 & 86\\
Gap candidate & 2,833 & 154,928 & 0\\
\bottomrule
\end{tabularx}
\caption{Static measurements of the final, audit-inclusive, uncompiled
candidate files. All 1,192 public theorem header strings were compared
byte-for-byte; this is not an elaborated theorem-type comparison.}
\end{table}

\textbf{The result is 660 fewer physical lines but 2,360 more bytes.}
It demonstrates less tactic repetition, not byte compression or a speedup.
Much of the line change removes proof-body layout. The cost of one finite
certificate could exceed that of the old local proofs; benchmark before
adoption, and split it into checked chunks if a monolithic check is costly.

There is an additional reuse opportunity, not a new data format: the Split
labels extracted for this review are exactly the integer list already stored
as \code{SplitJointPower59.natural_oldLabelsData} at
\src{SplitJointPower59.lean}{64}{78}. Extract that table and its certificate
into a small shared parent module rather than duplicate it or import the
joint child circuit merely to obtain parent metadata.

\subsection{Offer an indexed facade before removing named wrappers}
\label{sec:indexed}

\textbf{Priority 3.} Each gate wrapper specializes
\code{run_node_equation}. Instead of materializing every instance, expose
one theorem indexed by \code{Fin nodeCount}. The proposed Gap facade is:

\begin{Lean}
theorem natural_step (p : Fin 11 → R) (x : Fin 9 → R) (i : Fin 484) :
    naturalRegister p x i.val =
      (GapCircuit66.naturalNodes.getD i.val (.literal 0)).evalAt p x
        (naturalRegister p x) :=
  run_node_equation p x GapCircuit66.naturalNodes
    GapCircuit66.natural_topological i.val
    (by simpa only [GapCircuit66.natural_length] using i.isLt)
\end{Lean}

The supplied compact facades are 94 lines for Split and 92 for Gap. They
contain the register evaluator, indexed step theorem, complete label data,
power theorem, concrete examples, and axiom probes. \textbf{They are not
drop-in replacements}: their namespaces differ and the 1,192 old named
theorems are absent. Their size is not a measured project-wide saving.

Start with one genuine consumer such as
\code{GapCircuitRegisters66.natural_baseRegister} at
\src{GapCircuitRegisters66.lean}{18}{22}. Check that a concrete indexed
instance rewrites without unfolding the whole evaluator. Then retain useful
semantic boundary names such as \code{natural_kernel0Register},
\code{natural_sourceHighRegister}, and \code{natural_sourceRankRegister}
while gradually retiring only the mechanical wrappers.

The direct Split consumers are Registers59 and Translation59. The direct
Gap consumers are Degree66, Tail66, Translation66, and Registers66; recipes
depend on them transitively. Rebuild the reverse dependency closure, not
just these direct imports. In particular, the long, selective rewrite lists
at \src{GapCircuitRegisters66.lean}{30}{40} maintain deliberate cut points
in a shared expression graph. Replacing them with \code{simp [all definitions]}
can expand shared subexpressions repeatedly.

\subsection{Certify semantic blocks instead of globally normalizing the DAG}

\textbf{Priority 3, only after a consumer pilot.} The existing Gap Registers
file already establishes 13 semantic register facts at
\src{GapCircuitRegisters66.lean}{18}{94}, and the recipe composes them at
\src{GapCircuitRecipe66.lean}{15}{211}. This is the right proof structure.
Generalize it only where the kernel-block argument truly repeats.

A possible next layer is a block certificate with an input/base register,
sparse coefficient/exponent data, and an output register. A local
symbolic-value table can use \code{value_eq_of_node_equations} to establish
the block's meaning. It must also prove that the certified block embeds into
the \emph{existing full DAG} and has the claimed output register. Proving a
freshly generated, similar-looking circuit correct is not enough.

Keep the independently transcribed high-level recipe as the specification.
Defining the specification to be the same generated DAG would make agreement
vacuous. \code{SplitJointPower59} already demonstrates genuine certificate
reasoning between distinct literal circuits: the natural parent has 515
nodes and the joint child 498. Those are intentionally different objects.

\subsection{Share the ShortSystem degree prefix before changing representation}
\label{sec:shortdegree}

\textbf{Priority 3.} The two generated degree sections in
\code{ShortSystem} are at \src{ShortSystem.lean}{565}{842} and
\src{ShortSystem.lean}{844}{1115}. They contain respectively 96 and 94
local \code{let}/\code{have} pairs and share an exact 224-line prefix
beginning at the first local binding.

The least disruptive change is a common core-degree lemma exposing only
the boundary facts needed by the two output tails. A larger alternative uses
the existing \code{ArithmeticDAGDegree} checker for a shared literal prefix
and two residual/output tails. Either way, retain the bridges to
\code{literalValue}/\code{compactLiteralValue} and the original polynomial
values at \src{ShortSystem.lean}{348}{359} and
\src{ShortSystem.lean}{487}{498}.

The public bound is degree \emph{at most} 68, not exact degree 68. Preserve
that bound. A syntactic degree checker can miss cancellations; the explicit
warning at \src{GapCircuitDegree66.lean}{67}{67} documents this limitation
elsewhere in the same development. A weaker bound is not an acceptable
substitute merely because it is easier to certify.

\subsection{Treat typed circuit views as a later design option}

\textbf{Priority 4.} \code{ArithmeticDAG.Node} already uses \code{Fin} for
input and fixed-leaf indices, but natural numbers for operand registers
(\src{ArithmeticDAG.lean}{11}{18}). Invalid register access returns zero
(\src{ArithmeticDAG.lean}{79}{81}); a separate topology certificate excludes
invalid edges for well-formed circuits
(\src{ArithmeticDAG.lean}{105}{112}).

A small checked-circuit wrapper can bundle the existing list, its topology
proof, and a bounded output index. A dependent \code{SafeNode} view with
operands in \code{Fin earlier} is possible later. In both cases prove that
erasing the view returns \emph{exactly} the original literal list in its
original order. That equality is the bridge preserving operation counts,
liveness, data audits, and the existing semantics.

This does not verify a Python importer or prove specification equivalence
by itself. Dependent equalities can increase elaboration cost. It is lower
priority than the direct reuse of \code{run_power}, and no timing benefit is
claimed. In particular, none of the circuit proposals changes the original
natural operation counts, 480 for Split at
\src{SplitCircuit59.lean}{1246}{1246} and 450 for Gap at
\src{GapCircuit66.lean}{1186}{1186}.
